% Esquisse d'un programme (cote 119, copy 2): archive cote 119, site page 112 (PDF page 113), typescript page 34.
% typescript. Draft compared: esquisse/pages/esq_119_036.tex (copy 1, site 36).
% AI-transcribed by Claude Opus 5.5 (claude-opus-5-5, adaptive thinking, effort high), Message Batches API, batch esq119_full (msgbatch_015xqDHkTumDysRF1qngZhRJ), 2026-11-13; not yet checked by eye.
\scanpage{119/112}
\mspage{34}

le foncteur  $(U,Y,f) \longmapsto (X,Y)$   est défini de façon évidente,
indépendamment de toute théorie de voisinages tubulaire.

On dira sans doute, non sans quelque raison, que tout
celà n'est peut-être que rêves, qui s'évanouiront en fumée dès
qu'on s'essayera à un travail circonstancié, voire même dès avant
en face de certains faits connus ou bien évidents qui m'auraient
échappé. Certes, seul un travail sur pièces permettra de décan-%
ter le juste du faux et de connaître la substance véritable. La
seule chose dans tout cela qui ne fait pour moi l'objet d'aucun
doute, c'est la nécessité d'un tel travail de fondements, en
d'autres termes, la nature articielle des fondements actuels de
la topologie, et des difficultés que ceux-ci soulèvent à chaque
pas. Il est bien possible par contre que la formulation que je
donne à une théorie de dévissage des structures stratifiées, com-%
me un théorème d'équivalence de catégories isotopiques (voire même
$\infty$-isotopiques) convenables, soit trop optimiste. Je devrais ajou-%
ter pourtant que je n'ai guère de doutes non plus que la théorie
de ces dévissages que j'ai développé il y a deux ans, alors qu'elle
reste partiellement heuristique, exprime bel et bien une réalité
tout ce qu'il y a de palpable. Dans une partie de mon travail,
faute de pouvoir disposer d'un contexte ``modéré'' tout fait, et
pour avoir néanmoins des énoncés précis et démontrables, j'ai été
amené à postuler sur la structure stratifiée de départ des struc-%
tures supplémentaires tout ce qu'il y a de plausible, dans la na-%
ture de la donnée de rétractions locales notamment, qui dès lors
permettent bel et bien la construction d'un système canonique
d'espaces, paramétré par l'ensemble ordonné des ``drapeaux'' Drap(I)
de l'ensemble ordonné  $I$  indexant les strates, ces espaces jouant
le rôle des espaces  $(U,Y)$  de tantôt, reliés entre eux par des
applications de plongements et de fibrations propres, qui permet-%
tent de reconstituer de façon toute aussi canonique la structure
stratifiée de départ, y compris ces ``structures supplémentaires''(${}^{7}$).
Le seul ennui, c'est que ces dernières semblent un élément de struc-%
ture superfétatoire, qui n'est nullement une donnée dans les situations

\margin{112}
